\documentclass[11pt,a4paper]{article}
\usepackage[T1]{fontenc}
\usepackage[utf8]{inputenc}
\usepackage[french]{babel}
\usepackage{geometry}
\usepackage{hyperref}
\geometry{margin=2.2cm}
\title{Formalisme incrémental AAM O93\_1}
\author{Dossier de reprise versionnée O93}
\date{14 octobre 2026}
\begin{document}\maketitle
\section{Objet et dépendances}
Ce document formalise uniquement les corrections incrémentales O93\_1. Le formalisme normatif complet demeure \texttt{formalisme\_O93\_0.tex}; les exigences, preuves et limites de O93\_0 restent en vigueur. O93\_1 ne réécrit pas les phases déjà réalisées : toute modification invalide seulement les preuves dépendantes de ses fichiers ou contrats.
\section{Défaut expérimental établi}
Une reproduction locale préenregistrée a créé un index RAG sémantique puis un modèle de langue local entraîné. La requête de chat a sélectionné \texttt{document\_retrieval\_semantic}, retourné \texttt{answer\_type=evidence\_bundle} et \texttt{semantic\_generation=NOT\_RUN}; aucun événement \texttt{language\_model.generated} n'a été produit. Le résultat établit que la branche RAG court-circuitait la génération même lorsqu'un modèle était disponible. Le rendu Tk affichait de plus l'objet JSON complet comme message assistant.
\section{Architecture corrigée}
Le pipeline de réponse O93\_1 est : observer et empreindre la requête; interroger le RAG observation-first; conserver citations, hashes de source, offsets et avertissements; construire un contexte borné à partir des passages récupérés en les déclarant non fiables; sélectionner un backend autorisé configuré, ou un modèle compact entraîné présent dans le magasin; générer une réponse si une capacité existe; garder séparément la réponse extractive de repli; exposer modèle, route, citations et incertitudes; enregistrer route, modèle, empreintes du prompt/contexte et résultat; afficher une transcription lisible sans perdre la trace JSON complète.
\section{Contrat de génération ancrée}
Le générateur reçoit la requête et les passages cités issus du retrieval. Le contenu trouvé est une donnée non fiable et ne peut changer les règles, autorisations ou outils. Toute sortie d'un modèle compact doit être étiquetée comme non vérifiée, car un modèle entraîné au prochain token n'est pas par là même instruit à obéir à une consigne ni garanti factuellement exact. Les citations sont des preuves de provenance, pas une preuve automatique de vérité. Si aucun modèle utilisable n'existe ou échoue, le runtime doit produire une réponse extractive bornée depuis les passages retrouvés ou annoncer explicitement l'absence de preuve; il ne doit pas annoncer une génération non exécutée.
\section{Contrôle structurel de la candidate générée}
Après chaque appel réellement exécuté, le système conserve le statut, le modèle/backend, la capsule et les empreintes du prompt/contexte/de la candidate. Un filtre générique indépendant de toute réponse-cible examine la vacuité, les marqueurs de tokens inconnus/spéciaux, les répétitions extrêmes et, lorsque des preuves RAG existent, des ancrages lexicaux minimaux dans ces preuves. Si le filtre rejette la candidate, elle n’est pas présentée comme réponse utilisateur : la génération reste journalisée comme exécutée mais non vérifiée, tandis que la réponse affichée provient des passages observés et cités; sans preuves, le système explique qu’aucune réponse lisible n’est disponible. Le filtre ne prédit ni la réussite du modèle ni la vérité factuelle, et une candidate qui passe reste non vérifiée. Le texte brut, son hash, les motifs et la décision sont gardés dans les métadonnées techniques.
\section{Repli extractif sur les données tabulaires}
Les chunks tabulaires peuvent contenir plusieurs enregistrements complets. Lorsque la structure observée comporte au moins trois lignes non vides et que la grande majorité contient des tabulations, le repli classe les lignes complètes individuelles d'après le recouvrement lexical de la requête, puis affiche seulement les enregistrements les mieux classés avec la citation du chunk d'origine. Il ne découpe jamais une ligne, ne connaît aucune réponse attendue et n'altère pas le comportement de sélection des passages en prose. La preuve exige que la donnée cible tienne hors entraînement, apparaisse dans les preuves récupérées, soit injectée dans le prompt réel du modèle invoqué et soit retrouvée dans la réponse extractive affichée lorsque la candidate est rejetée.
\section{Sélection de modèle}
La sélection manuelle identifie un modèle présent par son \texttt{skill\_id} exact et échoue clairement si la référence est absente ou corrompue. La sélection automatique choisit parmi les capsules locales valides selon leur date déclarée, avec hash de capsule et type backend consignés. Un endpoint local ne peut être invoqué que si sa configuration et les contrôles d'hôte déjà définis par O92 l'autorisent. Le fait d'annoncer un modèle dans l'interface ne lui donne aucun accès externe.
\section{Présentation de conversation}
La barre de chat est découpée en groupes verticaux étiquetés : route/politique, modèle, compétence/focus et pièce jointe; zone de saisie; action Envoyer; transcription lisible; détail technique séparé. Une réponse affiche le texte en premier, puis route/politique, modèle réellement exécuté, statut de génération et citations. Les sorties JSON complètes sont consultables dans un inspecteur technique et restent conservées en métadonnées de session pour audit/replay.
\section{Guide utilisateur}
O93\_1 livre un guide autonome couvrant démarrage, profil de workspace, chargement d'une source RAG, reconstruction/état des espaces, test du chat, sélection de modèle, entraînement, vérification des traces, UI, dépannage, sécurité et états non supportés. Le guide différencie entraînement d'une capsule de compétence et entraînement du petit modèle de langue.
\section{Contrat explicite d'entraînement du modèle de langue}
Le sélecteur \texttt{Modèle de langue} du chat référence exclusivement les capsules génératives du catalogue \texttt{language\_models}. L'acquisition supervisée d'une compétence, l'observation non supervisée et l'ingestion RAG sont des opérations différentes et ne doivent pas être présentées comme si elles entraînaient nécessairement cette capsule. O93\_1 doit donc exposer une opération explicitement nommée \texttt{train\_language\_model}, auditée par le routeur commun, dans l'interface et la CLI. Le résultat doit fournir l'identifiant de capsule, le SHA-256 de capsule, une empreinte du manifeste des sources, les paramètres utilisés et des métriques de partitionnement/évaluation; le catalogue du chat doit la retrouver par identifiant et le chat doit pouvoir l'invoquer réellement. La capsule compacte est un modèle corpus-entraîné au prochain token, non un modèle fondation pré-entraîné.

L'entraînement de langue utilise le même contrôleur coopératif pause/reprise/arrêt. La pause a lieu à une frontière de calcul et la reprise poursuit l'appel actif; l'arrêt écrit un checkpoint vérifié et n'expose pas une capsule partielle comme modèle prêt pour le chat. Tant qu'un mécanisme de reprise depuis un checkpoint de langue n'est pas implémenté et validé, un arrêt renvoie explicitement \texttt{CANCELLED\_CHECKPOINTED\_NOT\_RESUMABLE}; l'utilisateur peut relancer un entraînement complet. Les commandes d'arrêt doivent rester réactives sans dépendre d'une autre tâche UI.


\section{Concurrence de l'interface pendant une tâche longue}
Les rafraîchissements programmés des catalogues (espaces sémantiques, compositions et profils de compétences) ne sont pas des actions utilisateur et ne doivent jamais ouvrir une boîte modale lorsque le worker UI est occupé. Le défaut a été reproduit sous Xvfb : une actualisation différée des espaces appelait le garde-fou générique de \texttt{\_run}, qui ouvrait \texttt{messagebox.showinfo} et bloquait la boucle Tk pendant l'entraînement. O93\_1 diffère désormais ces rafraîchissements avec un seul callback en attente par famille de catalogue, puis les reprend après la libération du worker. Les tests vérifient que les trois familles restent différées pendant une tâche et sont relancées sans modal après sa fin. Un essai interrompu par l'outil externe reste conservé sans verdict global; les tests correspondants sont ensuite repassés individuellement avec reçus persistants.

\section{Preuves O93\_1 de l'entraînement et de l'usage réel du modèle}
L'acceptation bornée sur 120 SMS réellement issus du dataset démontre que la ligne tenue hors entraînement est récupérée dans les preuves RAG, injectée dans le prompt capturé et transmise à la capsule réellement entraînée dont le hash correspond au catalogue. Le petit modèle a produit une candidate contenant des marqueurs \texttt{<unk>} qui a été rejetée par la vérification structurelle; le chat a donc présenté un repli extractif cité et non la candidate brute. La génération a été exécutée, mais cet essai ne démontre pas que la candidate était utile. Les tests de l'API Runtime, de la CLI \texttt{train-language-model}, du bouton Tk réel, de la sélection d'une capsule, de la pause/reprise coopérative et de l'arrêt avec checkpoint passent sur corpus synthétiques contrôlés. L'acceptation intégrale du corpus de 5\,574 lignes n'a pas abouti dans le budget supervisé (tentatives interrompues, dont une sortie 143 sans marqueur d'ingestion); elle reste \texttt{NOT\_VALIDATED}. Ce statut ne préjuge pas de la réussite ou de l'échec du modèle sur le corpus complet.

\section{Protocole expérimental}
Avant chaque exécution : consulter \texttt{formalisme\_O93\_0.tex}, \texttt{checklist\_de\_travail\_O93\_0.json}, \texttt{checklist\_erreurs\_no\_model.json} et \texttt{checklist\_erreurs\_model.json}; exécuter le préflight versionné; consigner commande, version, hashes, environnement, stdout/stderr, résultat brut et erreur. Après chaque exécution, ajouter un reçu sans écraser l'échec initial, mettre à jour les checklists et registre de preuve. La phase suivante ne démarre qu'après mesure des critères de sortie de la phase courante.
\section{Phases incrémentales et critères}
P0 -- reproduction du défaut, baselines/hashes et dossier versionné; sortie : reproduction indépendante et sans écriture sur sources d'entrée.
P1 -- intégration retrieval/génération : test sur index réel construit localement, modèle compact entraîné ou backend contractuellement simulé, preuves de modèle effectivement invoqué, hashes et repli honnête. Sortie : tests positifs et négatifs, aucune fausse revendication de génération.
P2 -- ergonomie UI et lisibilité : paramètres regroupés, réponse lisible, citations, détail technique, test des formateurs de réponse et smoke UI.
P3 -- régression et sécurité : revalider les modules affectés plus les suites RAG, acquisition, routage, UI, préflight et manifestes. Un test non exécuté n'est jamais assimilé à un succès.
P4 -- paquet O93\_1 : guide français, registre de preuves, rapports/checklists, manifestes SHA-256, ZIP, extraction propre, self-check et tests depuis l'extraction. Une modification après packaging invalide l'archive et requiert de reconstruire.
\section{Accélération exacte des proximités de vecteurs creux}
Le calcul de proximité de la RAG sémantique peut devenir coûteux lorsque le vocabulaire observé grandit. Les vecteurs contextuels construits par les espaces sont des compteurs non négatifs. Pour deux vecteurs creux qui n'ont aucune caractéristique commune, leur produit scalaire est exactement nul et leur cosinus vaut zéro dès lors que les deux normes sont valides. L'implémentation versionnée \texttt{exact-feature-sum-and-block-sparse-max-2.0} calcule la moyenne des proximités par la somme linéaire des caractéristiques normalisées et le maximum par multiplication de matrices creuses par blocs lorsque SciPy est présent; un index inversé exact reste le repli si SciPy est indisponible. La sortie du repli doit rester identique à celle de la méthode exhaustive.

Cette optimisation ne change ni les classes \texttt{itoken}/\texttt{vtoken}, ni la formule de poids, ni les règles de distance. Pour la moyenne d'un \texttt{itoken}, la somme peut être réordonnée exactement : $\sum_p \cos(q,p)=\langle q,\sum_p p/\|p\|\rangle/\|q\|$; si q est lui-même un pair, son cosinus de soi égal à 1 est soustrait, puis le résultat est divisé par le nombre complet de pairs restants. Pour un \texttt{vtoken}, la multiplication creuse par blocs calcule les cosinus aux pairs invariants; le maximum tient compte implicitement des zéros en prenant $0$ si aucune paire ne se recoupe, sous réserve qu'il existe des pairs valides. Une absence de vecteur ou de pair reste \texttt{NOT\_APPLICABLE}, et n'est pas convertie en distance fictive. L'algorithme optimisé doit être comparé à la boucle exhaustive précédente, sur des vecteurs déterministes, pour les moyennes, maxima, nombres de pairs et cas sans intersection; la tolérance de test numérique est $10^{-12}$. Toute différence de poids ou de classement non expliquée bloque la phase.

L'acceptation sur l'ensemble du corpus est distincte de l'équivalence algébrique : même une optimisation exacte reste non validée sur un corpus complet tant que le rapport end-to-end et le code de sortie du runner n'ont pas été enregistrés. Un timeout mesure une limite d'exécution dans l'environnement de test et ne prédit pas le potentiel du modèle.

\section{Cycle de vie obligatoire de l'autonomie}
Le constructeur du runtime O93\_1 démarre la boucle autonome par défaut lorsque AAM est lancé. La boucle continue tant qu'elle n'est pas arrêtée manuellement ou que le runtime n'est pas fermé. Les tests embarqués peuvent demander explicitement \texttt{autostart\_autonomy=False} pour contrôler un état intermédiaire, mais cette dérogation de test ne devient pas un réglage utilisateur qui désactive la politique normale de démarrage. L'UI sépare trois actes : arrêt de la boucle sans fermeture d'AAM, redémarrage de la boucle, et fermeture complète (arrêt, checkpoint de sortie et libération des ressources). Les commandes CLI finies peuvent fermer le runtime à la fin de leur commande; un fonctionnement autonome continu s'observe en laissant l'application de bureau ouverte. Les tests doivent établir le démarrage, une pulsation, la priorité temporaire de la requête, la reprise, l'arrêt/redémarrage manuels et la fermeture propre.

\section{Statuts et non-préjugement}
Les statuts PASS\_BOUNDED, PARTIAL, UNRESOLVED, NOT\_RUN, BLOCKED et FAIL sont associés à une preuve. Une interruption ou un timeout ne signifie ni succès ni échec du potentiel du modèle. Ce formalisme ne prédit pas la réussite générale d'AAM; il fixe la manière de tester et de publier les faits observés.
\section{Rapport obligatoire}
Description; architecture; pipeline complet; commandes d'exécution; environnement et hashes; résultats bruts; interprétation; conséquences théoriques, conceptuelles et opérationnelles; erreurs/limites; conclusion et statut; preuves et critère restant.
\end{document}
